% ============================================================
%  part04.tex  |  原文 PDF p.23–30（页脚 p.22–29）
%  内容：第 4 章 对称性的代数方面（Algebraic Aspects of Symmetry）
%        4.1 Chevalley--Eilenberg 复形（PDF p.23–24）
%        4.2 微分分次 Lie 代数（DGLA）（PDF p.25–28）
%        4.3 同伦 Lie 代数（$L_\infty$-代数）（PDF p.29–30）
% ============================================================

% >>> p023
\section{对称性的代数方面（Algebraic Aspects of Symmetry）}

\begin{keypoints}
  \item \textbf{核心物理图像}：规范理论的代数骨架是“物质场 $+$ 鬼场 $+$ BRST 算子”构成的
        上链复形；经典对称性由 Lie 代数（更一般地，DGLA）描述，量子化之后其代数关系
        只在“同伦”意义下成立，这正是 $L_\infty$ 代数的由来。
  \item \textbf{数学工具}：外代数 $\bigwedge^{\bullet}\mathfrak g^{\vee}$、分次对称/反对称张量、
        Chevalley--Eilenberg 微分 $\dd_{\mathrm{CE}}$、微分分次 Lie 代数（DGLA）及其
        Chevalley--Eilenberg 复形、$L_\infty$ 代数。
  \item \textbf{主要结论}：$\dd_{\mathrm{CE}}^2=0$ 等价于 Lie 括号的 Jacobi 恒等式；对 DGLA 有
        $\dd_{\mathrm{CE}}=\dd_1+\dd_2$，且 $\dd_{\mathrm{CE}}^2=0$ 等价于 DGLA 的定义公理；$L_\infty$
        关系式 $\sum\pm\ell(\ell)=0$ 等价于对偶侧 $\delta^2=0$，后者可以画成 Feynman 图恒等式。
  \item \textbf{译者评注}：本节是全书的“代数语法”：它把“对称性”翻译成“某个分次代数上的
        微分 $\dd$，且 $\dd^2=0$”。第 6 章的 BRST--BV 形式不过是把这里的
        $\bigwedge^{\bullet}\mathfrak g^{\vee}$ 换成场空间上的多向量场代数；4.3 节的
        $L_\infty$ 关系则将在第 5、7 章以 Feynman 图与同伦重整化群流的形式重现。
\end{keypoints}

\subsection{Chevalley--Eilenberg 复形（Chevalley--Eilenberg Complex）}

我们讨论规范理论背后的代数结构。设 $G$ 是 Lie 群，$V$ 是相应的 $G$-表示空间。
在规范理论中，$G$ 将是规范变换构成的无穷维（$\infty$-维）Lie 群，而 $V$ 代表物质场。

为简化表述并说明主要结构，本节我们把注意力集中在一个玩具模型（toy model）上：
$G$ 是有限维连通 Lie 群，其 Lie 代数为 $\mathfrak g$，$V$ 是有限维表示。令
\[
  \mathfrak g^{\vee}=\Hom(\mathfrak g,k)
\]
为 $\mathfrak g$ 的线性对偶（基域 $k=\R$ 或 $\C$）。考虑外代数
\[
  \bigwedge{}^{\bullet}\mathfrak g^{\vee}
  =\bigoplus_{p}\bigwedge{}^{p}\mathfrak g^{\vee}.
\]
设 $\{e_1,\dots,e_m\}$ 是 $\mathfrak g$ 的一组基，其中 $m=\dim\mathfrak g$，而
$\{c^1,\dots,c^m\}$ 是 $\mathfrak g^{\vee}$ 的对偶基。于是可以显式地写成
\[
  \bigwedge{}^{\bullet}\mathfrak g^{\vee}=k[c^\alpha],
  \qquad\text{其中}\quad c^\alpha c^\beta=-c^\beta c^\alpha .
\]
代数结构由楔积给出。设
\[
  [-,-]:\ \bigwedge{}^{2}\mathfrak g\longrightarrow\mathfrak g
\]
为 Lie 括号。其对偶给出映射
\[
  \dd_{\mathrm{CE}}:\ \mathfrak g^{\vee}\longrightarrow\bigwedge{}^{2}\mathfrak g^{\vee},
\]
它可唯一地延拓为映射
\[
  \dd_{\mathrm{CE}}:\ \bigwedge{}^{\bullet}\mathfrak g^{\vee}
  \longrightarrow\bigwedge{}^{\bullet}\mathfrak g^{\vee},
\]
由以下条件确定：
\begin{enumerate}[label=\arabic*),leftmargin=1.8em,itemsep=0.2em]
  \item 在生成元上，$\dd_{\mathrm{CE}}:\mathfrak g^{\vee}\to
        \bigwedge^{2}\mathfrak g^{\vee}$ 如上；
  \item 分次 Leibniz 法则
  \[
    \dd_{\mathrm{CE}}(a\wedge b)
    =(\dd_{\mathrm{CE}}a)\wedge b+(-1)^{p}\,a\wedge \dd_{\mathrm{CE}}b,
    \qquad\text{其中 } a\in\bigwedge{}^{p}\mathfrak g^{\vee}.
  \]
\end{enumerate}

% [Tutor Note: 分次符号规则是全节的“事故高发区”。此处 $(-1)^{p}$ 中的 $p$ 必须是
% 元素 $a$ 的**总次数**（本例中 $a\in\bigwedge^{p}\mathfrak g^{\vee}$ 是齐次元，故 $p=|a|$）；
% 若 $a$ 非齐次须先按次数拆开再逐项应用。符号的来源是 $\dd_{\mathrm{CE}}$ 为奇导子：
% 要把 $\dd_{\mathrm{CE}}$ 移过 $a$ 才能作用到 $b$ 上，每一次“移过”贡献一个 $-1$。
% 注意这个符号只与“搬动 $d$”有关，与 $d$ 本身作用在什么上无关。]

用上述基显式写出：
\[
  [e_\alpha,e_\beta]=\sum_{\gamma}f^{\gamma}_{\alpha\beta}\,e_\gamma,
\]
其中 $f^{\gamma}_{\alpha\beta}$ 是结构常数。对偶地，在生成元上有
\[
  \dd_{\mathrm{CE}}c^{\gamma}
  =\frac{1}{2}f^{\gamma}_{\alpha\beta}\,c^{\alpha}c^{\beta}.
\]
于是 $\dd_{\mathrm{CE}}$ 在一般元素上的作用为
\[
  \begin{aligned}
    \dd_{\mathrm{CE}}\bigl(c^{\alpha_1}c^{\alpha_2}\cdots c^{\alpha_p}\bigr)
    ={}&(\dd_{\mathrm{CE}}c^{\alpha_1})\cdots\gamma^{\alpha_p}
      -c^{\alpha_1}(\dd_{\mathrm{CE}}c^{\alpha_2})\cdots c^{\alpha_p}\\
      &+\cdots+(-1)^{p-1}c^{\alpha_1}\cdots c^{\alpha_{p-1}}
      \bigl(\dd_{\mathrm{CE}}c^{\alpha_p}\bigr).
  \end{aligned}
\]

% [Tutor Note: 原文疑误：上式右端第一项末尾的 $\gamma^{\alpha_p}$ 在全书中再无定义，
% 按 Leibniz 展开应为 $c^{\alpha_p}$；并且首项按展开应为
% $(\dd_{\mathrm{CE}}c^{\alpha_1})c^{\alpha_2}\cdots c^{\alpha_p}$（原文在省略号前漏掉了
% $c^{\alpha_2}\cdots$）。此处照录原文 $\gamma^{\alpha_p}$，提请核对原书 p.23 / PDF p.23。
% 展开式的符号规律是：第 $i$ 项前面带 $(-1)^{i-1}$，因为要把 $\dd_{\mathrm{CE}}$ 移过前 $i-1$ 个
% 奇数次数的 $c$。]

% [Tutor Note: $\dd_{\mathrm{CE}}c^{\gamma}=\tfrac12 f^{\gamma}_{\alpha\beta}c^{\alpha}c^{\beta}$
% 中的 $1/2$ 极易漏写或写成 $f^{\gamma}_{\beta\alpha}$。核对方式：用对偶配对
% $\langle \dd_{\mathrm{CE}}c^{\gamma},e_\alpha\wedge e_\beta\rangle
% =c^{\gamma}([e_\alpha,e_\beta])=f^{\gamma}_{\alpha\beta}$，而右端
% $\tfrac12 f^{\gamma}_{\rho\lambda}\langle c^{\rho}c^{\lambda},e_\alpha\wedge e_\beta\rangle
% =\tfrac12\bigl(f^{\gamma}_{\alpha\beta}-f^{\gamma}_{\beta\alpha}\bigr)
% =f^{\gamma}_{\alpha\beta}$（结构常数对下指标反对称）。]

% >>> p024
\begin{wprop}{4.1}
  $\dd_{\mathrm{CE}}:\bigwedge^{\bullet}\mathfrak g^{\vee}
  \to\bigwedge^{\bullet}\mathfrak g^{\vee}$ 满足
  \[
    \dd_{\mathrm{CE}}^2=0,
  \]
  这等价于 $[-,-]$ 的 Jacobi 恒等式。
\end{wprop}

\begin{proof}
  留作练习。
\end{proof}

% [Tutor Note: 原文把证明留作练习（“Proof: Exercise.”），但这一步是理解全节的钥匙：
% $\dd_{\mathrm{CE}}^2c^{\gamma}$ 恰好给出结构常数的 Jacobi 恒等式
% $f^{\gamma}_{\alpha[\beta}f^{\alpha}_{\rho\lambda]}=0$（对三个下指标全反对称化后为零）。
% 建议读者真的动手算一遍——第 6 章 BRST 算子的 $Q^2=0$ 与此是同一个计算，
% 只是把常数 $f$ 换成含场的结构函数。]

\begin{wdefn}{4.2}
  $C^{\bullet}(\mathfrak g):=\bigl(\bigwedge^{\bullet}\mathfrak g^{\vee},
  \dd_{\mathrm{CE}}\bigr)$ 定义一个上链复形，其中
  \[
    C^{p}(\mathfrak g)=\bigwedge{}^{p}\mathfrak g^{\vee},
    \qquad\text{且}\qquad
    \dd_{\mathrm{CE}}:\ C^{p}(\mathfrak g)\longrightarrow C^{p+1}(\mathfrak g).
  \]
  这称为 \textbf{Chevalley--Eilenberg 上链复形}。
\end{wdefn}

\begin{wrem}{4.3}
  在物理（规范理论）的术语中，$c^{\alpha}$ 称为\textbf{鬼场}（ghost field），
  $\dd_{\mathrm{CE}}$ 称为 \textbf{BRST 微分}。
\end{wrem}

我们可以把这一构造推广，以容纳表示。设 $V$ 是一个 $\mathfrak g$-模，带有
\[
  \rho:\ \mathfrak g\otimes V\longrightarrow V.
\]
考虑自由的 $\bigwedge^{\bullet}\mathfrak g^{\vee}$-模
$\bigwedge^{*}\mathfrak g^{\vee}\otimes V$。存在唯一映射
\[
  \dd_{\mathrm{CE}}:\ \bigwedge^{p}\mathfrak g^{\vee}\otimes V
  \longrightarrow\bigwedge^{p+1}\mathfrak g^{\vee}\otimes V,
\]
由以下条件确定：
\begin{enumerate}[label=\arabic*),leftmargin=1.8em,itemsep=0.2em]
  \item $\dd_{\mathrm{CE}}:V\to\mathfrak g^{\vee}\otimes V$ 是表示映射
        $\rho:\mathfrak g\otimes V\to V$ 的对偶；
  \item 对 $a\in\bigwedge^{p}\mathfrak g^{\vee}$、$m\in V$，
  \[
    \dd_{\mathrm{CE}}(a\otimes m)
    =\dd_{\mathrm{CE}}(a)\otimes m+(-1)^{p}\,a\wedge \dd_{\mathrm{CE}}(m).
  \]
\end{enumerate}

\begin{wprop}{4.4}
  $\dd_{\mathrm{CE}}:\bigwedge^{\bullet}\mathfrak g^{\vee}\otimes V
  \to\bigwedge^{\bullet}\mathfrak g^{\vee}\otimes V$ 满足
  \[
    \dd_{\mathrm{CE}}^2=0.
  \]
  这等价于 $V$ 是一个 $\mathfrak g$-表示。
\end{wprop}

\begin{proof}
  留作练习。
\end{proof}

\begin{wdefn}{4.5}
  与 $\mathfrak g$-模 $V$ 相联系的 Chevalley--Eilenberg 上链复形是
  \[
    C^{\bullet}(\mathfrak g,V)
    :=\bigl(\bigwedge^{\bullet}\mathfrak g^{\vee}\otimes V,\,
    \dd_{\mathrm{CE}}\bigr).
  \]
\end{wdefn}

注意 $C^{\bullet}(\mathfrak g,V)$ 是一个 $C^{\bullet}(\mathfrak g)$-模：
\[
  C^{p}(\mathfrak g)\otimes C^{q}(\mathfrak g,V)\longrightarrow
  C^{p+q}(\mathfrak g,V),
  \qquad a\otimes u\longmapsto a\wedge u,
\]
它与微分相容：
\[
  \dd_{\mathrm{CE}}(a\wedge u)
  =(\dd_{\mathrm{CE}}a)\wedge u+(-1)^{p}\,a\wedge \dd_{\mathrm{CE}}(u).
\]
换言之，$C^{\bullet}(\mathfrak g,V)$ 是 $C^{\bullet}(\mathfrak g)$ 上的 dg-模。

% >>> p025
\subsection{微分分次 Lie 代数（Differential Graded Lie Algebra, DGLA）}

我们将在分次向量空间的范畴中工作。设
\[
  W=\bigoplus_{m\in\Z}W_m
\]
为 $\Z$-分次向量空间，其中 $W_m$ 是次数为 $m$ 的分量。对 $a\in W_m$，记 $|a|=m$。
我们将频繁使用分次空间上的以下运算。
\begin{itemize}[leftmargin=1.6em,itemsep=0.3em]
  \item $W[n]$ 表示 $\Z$-分次空间，其定义为 $W[n]_m:=W_{n+m}$（次数平移 $n$）。

\begin{figure}[htbp]
  \centering
  \includegraphics[width=0.74\linewidth]{fig_p025_1.png}
  \caption*{次数平移 $W\to W[n]$ 的示意图（原文 PDF p.25 示意图，原文无编号）}
\end{figure}

  \item $W^{\vee}$ 表示线性对偶，其中
        \[
          (W^{\vee})_m:=\Hom(W_{-m},k)=(W_{-m})^{\vee}.
        \]
  \item 给定两个 $\Z$-分次向量空间 $V,W$：
        \[
          (V\otimes W)_n:=\bigoplus_{i+j=n}(V_i\otimes W_j),
          \qquad
          \Hom(V,W)_n:=\bigoplus_{i}\Hom(V_i,W_{i+n}).
        \]
  \item $\Sym^{m}(V):=V^{\otimes m}/\sim$，其中
        \[
          a\otimes b\sim(-1)^{|a||b|}\,b\otimes a.
        \]
        这称为 $m$ 次\textbf{分次对称张量}。我们记
        \[
          \Sym(V):=\bigoplus_{m\ge0}\Sym^{m}(V),
          \qquad
          \widehat{\Sym}(V):=\prod_{m\ge0}\Sym^{m}(V).
        \]
  \item $\bigwedge^{m}(V):=V^{\otimes m}/\sim$，其中
        \[
          a\otimes b\sim-(-1)^{|a||b|}\,b\otimes a.
        \]
        这称为 $m$ 次\textbf{分次反对称张量}。
\end{itemize}

% [Tutor Note: 记号提醒：$W[n]_m:=W_{n+m}$ 意味着 $W[n]$ 中次数为 $m$ 的元素，原来次数为
% $m+n$，即 shift 把次数整体**降低** $n$。不同文献对 $[n]$ 的约定可能整体差一个符号，
% 引用公式时务必先确认。这里 $[-1]$ 的作用是把“鬼场”次数降 $1$，从而让
% $\mathfrak g^{\vee}[-1]$ 与通常外代数的次数对齐（见命题 4.6）。]

% [Tutor Note: 注意 $\Sym(V)$ 是直和、$\widehat{\Sym}(V)$ 是直积：前者只含有限个非零分量，
% 后者是形式幂级数（可含无穷多项）。原书 p.29 定义 $L_\infty$ 的导子 $\delta$ 时必须用
% 完备化的 $\widehat{\Sym}$，否则无穷和 $\delta=\sum_{n\ge1}\delta_n$ 无意义——
% 这正是原文括号中“so that $\delta$ is defined”的含义。原文 p.25 与 p.29 分别在
% $\Sym$ 上方加宽帽与波浪号，记号不统一，本译本统一写作 $\widehat{\Sym}$。]

\begin{wprop}{4.6}
  $\bigwedge^{k}(V[1])\cong\Sym^{k}(V)[k]$。
\end{wprop}

\begin{proof}
  这是一道很有用也很有益的练习。
\end{proof}

\begin{wdefn}{4.7}
  \textbf{分次 Lie 代数}是一个 $\Z$-分次空间
  \[
    \mathfrak g=\bigoplus_{m\in\Z}\mathfrak g_m,
  \]
  连同双线性映射
  \[
    [-,-]:\ \mathfrak g\otimes\mathfrak g\longrightarrow\mathfrak g,
  \]
  满足以下条件：
  \begin{enumerate}[label=\arabic*),leftmargin=1.8em,itemsep=0.25em]
% >>> p026
    \item 分次括号：$[\mathfrak g_\alpha,\mathfrak g_\beta]
          \subset\mathfrak g_{\alpha+\beta}$；
    \item 分次反对称性：
          \[
            [a,b]=-(-1)^{\alpha\beta}[b,a],
            \qquad\forall a\in\mathfrak g_\alpha,\ b\in\mathfrak g_\beta;
          \]
    \item 分次 Jacobi 恒等式：
          \[
            [[a,b],c]=[a,[b,c]]-(-1)^{\alpha\beta}[b,[a,c]],
            \qquad\forall a\in\mathfrak g_\alpha,\ b\in\mathfrak g_\beta,\
            c\in\mathfrak g_\gamma.
          \]
  \end{enumerate}
\end{wdefn}

% [Tutor Note: 分次 Jacobi 恒等式的符号 $(-1)^{\alpha\beta}$ 来自把次数为 $\alpha$ 的 $a$
% 移过次数为 $\beta$ 的 $b$。三条公理中符号最易记错的是这一条：若漏掉 $(-1)^{\alpha\beta}$
% 或把它写成 $(-1)^{\beta\gamma}$，$\dd_{\mathrm{CE}}^2=0$ 立刻失效。此外要注意原文的配对方式
% 是 $[a,[b,c]]$ 与 $[b,[a,c]]$，一旦改写成别的字母顺序，指数也要相应改变，务必重新推导
% 而不要硬背。]

\begin{wdefn}{4.8}
  \textbf{微分分次 Lie 代数}（DGLA）是一个分次 Lie 代数 $\mathfrak g$，连同
  一个次数为 $1$ 的线性映射 $\dd$（即 $\dd:\mathfrak g_m\to\mathfrak g_{m+1}$），满足
  \begin{enumerate}[label=\arabic*),leftmargin=1.8em,itemsep=0.25em]
    \item $\dd^2=0$，故 $(\mathfrak g,\dd)$ 是一个上链复形；
    \item 分次 Leibniz 法则：
          \[
            \dd[a,b]=[\dd a,b]+(-1)^{\alpha}[a,\dd b],
            \qquad\forall a\in\mathfrak g_\alpha,\ b\in\mathfrak g_\beta.
          \]
  \end{enumerate}
\end{wdefn}

\begin{wexam}{4.9}
  通常的 Lie 代数就是如下的 DGLA：
  \begin{itemize}[leftmargin=1.6em,itemsep=0.2em]
    \item $\mathfrak g=\mathfrak g_0$ 集中在次数 $=0$ 处；
    \item $\dd=0$。
  \end{itemize}
  因此 DGLA 是 Lie 代数的自然推广。
\end{wexam}

\begin{wexam}{4.10}
  设 $X$ 是一个流形，$\mathfrak g$ 是一个 Lie 代数。设 $(\Omega^{\bullet}(x),\dd)$ 是
  de Rham 复形。则 $(\Omega^{\bullet}(x)\otimes\mathfrak g,\,\dd,\,[-,-])$ 是一个 DGLA。
  \begin{itemize}[leftmargin=1.6em,itemsep=0.2em]
    \item $\Omega^{p}(x)\otimes\mathfrak g$ 是次数为 $p$ 的分量。
    \item $\dd:\Omega^{p}(x)\otimes\mathfrak g\to\Omega^{p+1}(x)\otimes\mathfrak g$ 是
          de Rham 微分：
          \[
            \dd(\alpha\otimes h)=\dd\alpha\otimes h,
            \qquad\text{对 } \alpha\in\Omega^{p},\ h\in\mathfrak g.
          \]
    \item 括号由 $\mathfrak g$ 上的 Lie 括号 $[-,-]_{\mathfrak g}$ 诱导：
          \[
            [\alpha_1\otimes h_1,\alpha_2\otimes h_2]
            :=(\alpha_1\wedge\alpha_2)\otimes[h_1,h_2]_{\mathfrak g},
            \qquad\forall\alpha_1,\alpha_2\in\Omega^{*},\ h_1,h_2\in\mathfrak g.
          \]
  \end{itemize}
  这个例子与 Chern--Simons 理论相关。
\end{wexam}

\begin{wexam}{4.11}
  设 $X$ 是一个复流形。设 $(\Omega^{0,\bullet},\bar\partial)$ 是 Dolbeault 复形。
  用 $T^{1,0}_X$ 表示 $(1,0)$-向量场丛。则
  \[
    \bigl(\Omega^{0,\bullet}\bigl(X,T^{1,0}_X\bigr),\,
    \bar\partial,\,[-,-]\bigr)
  \]
  是一个 DGLA。具体地，设 $z^i$ 是局部全纯坐标。元素
  $\alpha\in\Omega^{0,p}\bigl(X,T^{1,0}_X\bigr)$ 可写成
  \[
    \alpha=\sum_{i,J}\alpha^{i}_{J}\,\dd\bar z^{J}\otimes
    \frac{\partial}{\partial z^{i}}.
  \]
  这里 $J=\{j_1<\cdots<j_p\}$ 是一个多重指标，且
  $\dd\bar z^{J}=\dd\bar z^{j_1}\wedge\cdots\wedge \dd\bar z^{j_p}$。于是微分为
  \[
    \bar\partial\alpha
    =\sum\bar\partial\bigl(\alpha^{i}_{J}\bigr)\wedge \dd\bar z^{J}\otimes
    \frac{\partial}{\partial z^{i}}
    =\sum_{l}\frac{\partial\alpha^{i}_{J}}{\partial\bar z^{l}}\,
    \dd\bar z^{l}\wedge \dd\bar z^{J}\otimes\frac{\partial}{\partial z^{i}}.
  \]
\end{wexam}

% >>> p027
给定两个元素
\[
  \alpha=\sum\alpha^{i}_{J}\,\dd\bar z^{J}\otimes\frac{\partial}{\partial z^{i}},
  \qquad
  \beta=\sum\beta^{i}_{M}\,\dd\bar z^{M}\,\frac{\partial}{\partial z^{i}},
\]
它们的括号由下式给出：
\[
  [\alpha,\beta]
  =\sum\Bigl(\alpha^{i}_{J}\frac{\partial\beta^{i}_{M}}{\partial z^{j}}
  -\beta^{i}_{M}\frac{\partial\alpha^{i}_{J}}{\partial z^{j}}\Bigr)
  \dd\bar z^{J}\wedge \dd\bar z^{M}\otimes\frac{\partial}{\partial z^{i}}.
\]
在次数 $=0$ 的分量上，这恰好是 $(1,0)$-向量场上通常的 Lie 括号。

% [Tutor Note: 原文疑误（指标错位）：括号公式中 $\alpha$ 与 $\beta$ 的上指标都印成了 $i$，
% 而求导变量是 $z^{j}$。按定义应为
% $\alpha^{j}_{J}\partial_j\beta^{i}_{M}-\beta^{j}_{M}\partial_j\alpha^{i}_{J}$
% （上指标 $j$ 与求导变量缩并、$i$ 是向量场指标，$\sum_{i,J,M}$ 之外还隐含对 $j$ 求和）。
% 此外 $\beta$ 的展开式原文漏印了 $\otimes$（印成
% $\beta^{i}_{M}\dd\bar z^{M}\frac{\partial}{\partial z^{i}}$），宜读作
% $\beta^{i}_{M}\dd\bar z^{M}\otimes\frac{\partial}{\partial z^{i}}$。均照录原文，请核对
% 原书 p.26--27 / PDF p.26--27。同类排印不一致还见于例 4.10：流形先记 $X$，其后公式里
% 却一律写作小写 $x$（$\Omega^{\bullet}(x)$），本译本照录原文，读者应读作 $X$。]

这个 DGLA 与复结构的形变有关。在物理上，它与 B-扭曲拓扑弦场论
（Kodaira--Spencer 引力）相关。

Chevalley--Eilenberg 复形的构造可以自然地推广到 DGLA。设
$(\mathfrak g,\dd,[-,-])$ 是一个 DGLA。定义分次空间
\[
  C^{\bullet}(\mathfrak g):=\Sym\bigl(\mathfrak g^{\vee}[-1]\bigr)
  =\bigoplus_{m}\Sym^{m}\bigl(\mathfrak g^{\vee}[-1]\bigr).
\]
由于 $\mathfrak g^{\vee}[-1]=(\mathfrak g[1])^{\vee}$，我们可以把
$C^{\bullet}(\mathfrak g)$ 看作 $\mathfrak g[1]$ 上的（多项式）函数。当 $\mathfrak g$
是通常的 Lie 代数时，命题 4.6 蕴含
\[
  C^{k}(\mathfrak g)=\Sym^{k}\bigl(\mathfrak g^{\vee}[-1]\bigr)
  \simeq\bigwedge{}^{k}\mathfrak g^{\vee}[-k],
\]
即 $\bigwedge^{k}\mathfrak g^{\vee}$ 位于次数 $k$ 处。令
\begin{itemize}[leftmargin=1.6em,itemsep=0.3em]
  \item $\dd:\mathfrak g\to\mathfrak g$ 的对偶为
        \[
          \dd_1:\ \mathfrak g^{\vee}[-1]\longrightarrow\mathfrak g^{\vee}[-1];
        \]
  \item 括号 $[-,-]:\bigwedge^{2}\mathfrak g\to\mathfrak g$ 的对偶为
        \[
          \dd_2:\ \mathfrak g^{\vee}[-1]\longrightarrow
          \Sym^{2}\bigl(\mathfrak g^{\vee}[-1]\bigr)
          \simeq\bigwedge{}^{2}\mathfrak g^{\vee}[-2].
        \]
\end{itemize}
注意 $\dd_1$ 与 $\dd_2$ 的次数均为 $1$。由于 $C^{\bullet}(\mathfrak g)$ 由
$\mathfrak g^{\vee}[-1]$ 自由生成，我们可以按以下方式把 $\dd_1$ 与 $\dd_2$ 延拓到
$C^{\bullet}(\mathfrak g)$ 上：
\begin{itemize}[leftmargin=1.6em,itemsep=0.3em]
  \item 在生成元 $\mathfrak g^{\vee}[-1]$ 上，按上述定义；
  \item 分次 Leibniz 法则：对 $\dd=\dd_1$ 或 $\dd_2$，
        \[
          \dd(a\wedge b)=(\dd a)\wedge b+(-1)^{|a|}\,a\wedge \dd b.
        \]
\end{itemize}
定义 \textbf{Chevalley--Eilenberg 微分}
\[
  \dd_{\mathrm{CE}}=\dd_1+\dd_2:\ C^{\bullet}(\mathfrak g)\longrightarrow
  C^{\bullet}(\mathfrak g).
\]

% [Tutor Note: 为什么对 DGLA 的 CE 复形要用 $\Sym(\mathfrak g^{\vee}[-1])$ 而不是
% $\bigwedge^{\bullet}\mathfrak g^{\vee}$？对集中在 $0$ 次的通常 Lie 代数，两者由命题 4.6
% 同构（$k$ 次部分都落在次数 $k$）；但对一般 DGLA，$\mathfrak g$ 自身带奇偶，
% 只有 $\Sym$ 才是“$\mathfrak g[1]$ 上的函数环”，$\bigwedge^{\bullet}$ 会给出错误的符号。
% 这一点在第 6 章 BV 形式中至关重要：那里 $C^{\bullet}(\mathfrak g)$ 就是多向量场代数。
% 另注：上式 $\dd(a\wedge b)$ 中的“$\wedge$”是原文写法，在对称代数里其实是乘法
% （可写成 $\dd(ab)$），原文记号混用；照着读即可。]

\begin{wprop}{4.12}
  $\dd_{\mathrm{CE}}^2=0$。复形 $(C^{\bullet}(\mathfrak g),\dd_{\mathrm{CE}})$ 也称为
  Chevalley--Eilenberg 上链复形。
\end{wprop}

\begin{proof}[证明梗概]
  我们说明为何这是对的，细节留给读者。事实上，如果我们把
  $\dd_{\mathrm{CE}}$ 表示为
\begin{figure}[htbp]
  \centering
  \includegraphics[width=0.35\linewidth]{fig_p027_1.png}
  \caption*{$\dd_{\mathrm{CE}}$ 的图形表示（原文 PDF p.27 示意图，原文无编号）}
\end{figure}
% >>> p028
  那么我们有
\begin{figure}[htbp]
  \centering
  \includegraphics[width=0.71\linewidth]{fig_p028_1.png}
  \caption*{$\dd_{\mathrm{CE}}^2=0$ 的图形展开（原文 PDF p.28 示意图，原文无编号）}
\end{figure}
  我们可以“看出”，$\dd_{\mathrm{CE}}^2=0$ 等价于 DGLA 的定义性质：
  \begin{itemize}[leftmargin=1.6em,itemsep=0.2em]
    \item $\dd^2=0$，
    \item Leibniz 法则：$\dd[-,-]=[\dd(-),-]\pm[-,\dd(-)]$，
    \item Jacobi 恒等式。
  \end{itemize}
\end{proof}

% [Tutor Note: 图形证明里的“$\sum$ permutations”一项不能漏：$\dd_{\mathrm{CE}}^2$ 的四类图
% 分别对应 DGLA 的三条公理——直线 $\dd^2=0$；$\dd$ 与顶点的两种相对位置给出分次 Leibniz
% 法则（差一个 Koszul 符号）；两个顶点复合的置换和给出 Jacobi 恒等式。把图翻译成公式时，
% 每个三价顶点贡献一个括号、每条标 $\dd$ 的线贡献一次 $\dd$，而“外线的排列”贡献符号。
% 这套“图$\,=\,$公式”的翻译规则在第 5 章的 Feynman 图展开中会系统化。]

推广到 $\mathfrak g$-模是直接的。

\begin{wdefn}{4.13}
  设 $\mathfrak g$ 是一个 DGLA。一个 \textbf{$\mathfrak g$-模}是一个上链复形
  $(M,\dd_M)$，连同双线性映射 $\mathfrak g\otimes M\xrightarrow{\ \cdot\ }M$，满足
  \begin{itemize}[leftmargin=1.6em,itemsep=0.25em]
    \item $\mathfrak g_\alpha\otimes M_\beta\longrightarrow M_{\alpha+\beta}$；
    \item $a\cdot(b\cdot m)-(-1)^{\alpha\beta}\,b\cdot(a\cdot m)=[a,b]\cdot m$，
          $\forall a\in\mathfrak g_\alpha,\ b\in\mathfrak g_\beta,\ m\in M$；
    \item $\dd_M(a\cdot m)=(\dd a)\cdot m+(-1)^{\alpha}\,a\cdot \dd_M m$，
          $\forall a\in\mathfrak g_\alpha,\ m\in M$。
  \end{itemize}
\end{wdefn}

等价地，让我们定义分次空间
\[
  \mathfrak g_M:=\mathfrak g\oplus M,
\]
并赋予微分 $\dd\oplus \dd_M$ 与括号 $[-,-]$，其中
\begin{itemize}[leftmargin=1.6em,itemsep=0.2em]
  \item 在 $\mathfrak g\otimes\mathfrak g$ 上，$[-,-]$ 是 $\mathfrak g$ 上的括号：
        $[-,-]:\mathfrak g\otimes\mathfrak g\to\mathfrak g$；
  \item 在 $\mathfrak g\otimes M$ 上，$[-,-]$ 是 $\mathfrak g$-模作用：
        $[-,-]:\mathfrak g\otimes M\to M$；
  \item 在 $M\otimes M$ 上，$[-,-]=0$。
\end{itemize}
于是
\[
  M\ \text{是}\ \mathfrak g\text{-模}\iff\mathfrak g_M\ \text{是 DGLA}.
\]
相应的 Chevalley--Eilenberg 复形是
\[
  C^{\bullet}(\mathfrak g,M)
  :=\bigl(\Sym(\mathfrak g^{\vee}[-1])\otimes M,\,\dd_{\mathrm{CE}}\bigr).
\]
$\dd_{\mathrm{CE}}$ 的描述我们留给读者。

% [Tutor Note: 第二组条件（$M\otimes M$ 上的括号取 $0$）说明 $\mathfrak g_M$ 是“半直积”
% Lie 代数 $\mathfrak g\ltimes M$ 的分次版本；而 $M\otimes M$ 上的括号为零，正是
% “$M$ 是模”而非“$M$ 是子代数”的代数表述。第一条式子
% $a\cdot(b\cdot m)-(-1)^{\alpha\beta}b\cdot(a\cdot m)=[a,b]\cdot m$ 就是分次 Jacobi
% 恒等式在 $\mathfrak g\oplus M$ 上的投影，写成分次形式时那个 $(-1)^{\alpha\beta}$ 极易漏掉。]

% >>> p029
\subsection{同伦 Lie 代数（Homotopy Lie Algebra，$L_\infty$-代数）}

给定一个分次向量空间 $V$，考虑 $\Sym(V)$ 上次数为 $1$ 的导子（derivation）
\[
  \delta:\ \Sym(V)\longrightarrow\Sym(V),
\]
它满足分次 Leibniz 法则
\[
  \delta(a\otimes b)=(\delta a)\otimes b+(-1)^{|a|}\,a\otimes\delta b.
\]
这样的 $\delta$ 完全由它在生成元上的作用决定：
\[
  \delta:\ V\longrightarrow\Sym(V).
\]
我们可以分解
\[
  \delta=\delta_0+\delta_1+\delta_2+\cdots,
\]
其中 $\delta_k:V\to\Sym^{k}(V)$。对 DGLA，我们有
\[
  \delta=\dd_{\mathrm{CE}}=\dd_1+\dd_2=\delta_1+\delta_2.
\]
这是一个只有 $\delta_1,\delta_2$ 非平凡的导子。把所有可能的分量 $\delta_k$ 都编码进来，
自然地推广这一算子，就引出了 $L_\infty$-代数。

% [Tutor Note: 原文从 $\delta_0$ 开始分解，但紧接着又只要求 $\delta_k:V\to\Sym^{k}(V)$，
% 而 $\Sym^{0}(V)=k$，故 $\delta_0$ 是一个次数 $1$ 的线性泛函 $V\to k$（仅当 $V$ 有
% $-1$ 次分量时可能非零）。通常的 $L_\infty$ 展开从 $\delta_1$ 起；$\delta_0\ne0$ 的情形
% 称为“弯曲”的（curved）$L_\infty$ 代数，本课程用不到。对 DGLA 有 $\delta_0=0$，
% 故原文随即写出 $\delta=\dd_1+\dd_2=\delta_1+\delta_2$。]

\begin{wdefn}{4.14}
  一个 \textbf{$L_\infty$-代数}是一个 $\Z$-分次向量空间 $\mathfrak g$，连同一族多重线性映射
  \[
    \ell_n:\ \bigwedge{}^{n}\mathfrak g\longrightarrow\mathfrak g,
    \qquad \deg(\ell_n)=2-n\quad(n\ge1),
  \]
  满足以下的 \textbf{$L_\infty$-关系式}
  \[
    \sum_{k=1}^{n}\pm\,\ell_{n-k+1}\bigl(\ell_k(\cdots),\cdots\bigr)=0,
    \qquad\forall n.
  \]
\end{wdefn}

% [Tutor Note: $\deg(\ell_n)=2-n$ 的正确性可用两个低阶项校验：$\ell_1$ 的次数为
% $2-1=1$，即微分（$\dd$）；$\ell_2$ 的次数为 $2-2=0$，即通常的 Lie 括号；$\ell_3,\ell_4,\dots$
% 是“高阶括号”，它们使得 Jacobi 恒等式只在同伦意义下成立。若把次数记成 $1-n$ 或
% $n-2$，则对偶侧导子 $\delta_n$ 的次数就不再是 $1$，整个构造失效。]

这种复杂的 $L_\infty$-关系可以这样理解。令
\[
  \delta_n:\ \mathfrak g^{\vee}[-1]\longrightarrow
  \Sym^{n}\bigl(\mathfrak g^{\vee}[-1]\bigr)
  \simeq\bigwedge{}^{n}\mathfrak g^{\vee}[-n]
\]
表示 $\ell_n:\bigwedge^{n}\mathfrak g\to\mathfrak g$ 的对偶。注意
\[
  \deg(\ell_n)=2-n\iff\deg(\delta_n)=1.
\]
令 $\delta=\sum_{n\ge1}\delta_n=\delta_1+\delta_2+\cdots$，它通过分次 Leibniz 法则定义了
$C^{\bullet}(\mathfrak g)=\widehat{\Sym}\bigl(\mathfrak g^{\vee}[-1]\bigr)$ 上的一个导子。
（这里我们使用（分次的）形式幂级数环
$\widehat{\Sym}\bigl(\mathfrak g^{\vee}[-1]\bigr)$，以保证 $\delta$ 有定义。）于是
\[
  \{\ell_n\}_{n\ge1}\ \text{的 } L_\infty\text{-关系}
  \iff \delta^2=0 .
\]
$(C^{\bullet}(\mathfrak g),\delta)$ 是 $\mathfrak g$ 的 Chevalley--Eilenberg 复形。

如果对每个 $\delta_n$ 都画一张图
\begin{figure}[htbp]
  \centering
  \includegraphics[width=0.24\linewidth]{fig_p029_1.png}
  \caption*{$\delta_n$ 的图形表示（原文 PDF p.29 示意图，原文无编号）}
\end{figure}
% >>> p030
那么 $\delta^2=0$ 可以画成
\begin{figure}[htbp]
  \centering
  \includegraphics[width=0.49\linewidth]{fig_p030_1.png}
  \caption*{$\delta^2=0$ 的图形表示（原文 PDF p.30 示意图，原文无编号）}
\end{figure}
正如我们将看到的，这通过 Feynman 图技术有一个自然的解释。关于 $L_\infty$-代数及其与
QFT 的关系，一个很好的综述见文献 [27] 及其参考文献。

% [Tutor Note: $L_\infty$-关系式里的“$\pm$”不是可有可无的装饰：它是 Koszul 符号，
% 由把次数为 $2-k$ 的映射 $\ell_k$ 移过其余输入所产生。$n=2$ 时该式化为
% $\ell_1\ell_2=\ell_2(\ell_1\otimes1\pm1\otimes\ell_1)$，即“$\ell_1$ 是 $\ell_2$ 的导子”
% （分次 Leibniz 法则）；$n=3$ 时它给出 Jacobi 恒等式的同伦修正：
% $\ell_1\ell_3+\ell_3\ell_1$ 与 $\ell_2\ell_2$ 只差一个 Koszul 符号的轮换和。
% 直接照抄别人的符号表是危险的——务必按“映射次数 + 输入置换”自行重算。]

% [Tutor Note: 物理图像点拨：$\delta^2=0$ 的图恒等式正是量子场论中 Ward 恒等式/
% 主方程的雏形——一个 $\ell_n$ 顶点代表 $n$ 个场的一次相互作用，$\delta^2=0$ 说
% “把所有可能的收缩方式加起来为零”。第 5 章的 Wick 定理与 Feynman 图展开、第 6--7 章的
% BV 主方程都会以不同的语言重复这张图。原文引用的 [27] 见书末参考文献。]
